\documentclass[11pt]{article}
\usepackage[a4paper,margin=1.9cm]{geometry}
\usepackage[T1]{fontenc}
\usepackage{textcomp}
\usepackage{amsmath,amssymb}
\usepackage{enumitem}
\usepackage{tikz}
\usetikzlibrary{calc,arrows.meta,patterns,decorations.pathmorphing}
\usepackage{xcolor}
\usepackage{multicol}
\usepackage{parskip}
\usepackage{lmodern}
\renewcommand{\familydefault}{\sfdefault}

\definecolor{unalblue}{RGB}{31,73,124}

\newlist{opciones}{enumerate}{1}
\setlist[opciones]{label=(\alph*),itemsep=1pt,topsep=2pt,leftmargin=1.8em,font=\normalfont}

\newcommand{\vect}[2]{\begin{pmatrix}#1\\#2\end{pmatrix}}
\newcommand{\vectiii}[3]{\begin{pmatrix}#1\\#2\\#3\end{pmatrix}}

\newcommand{\examheader}{%
  \noindent
  \begin{minipage}[t]{0.62\linewidth}
    {\bfseries\large Universidad Nacional de Colombia, sede Medell\'in}\\
    {\small Escuela de Matem\'aticas}\\
    {\small Primer Parcial de C\'alculo en Varias Variables}\\
    {\small 11 de marzo de 2017}
  \end{minipage}%
  \hfill
  \begin{minipage}[t]{0.33\linewidth}
    \raggedleft\small
    Nombre y c\'edula: \rule{2.8cm}{0.4pt}\\[4pt]
    Grupo: \rule{1.5cm}{0.4pt} \quad N\'umero asignado: \rule{1.5cm}{0.4pt}
  \end{minipage}\\[4pt]
  \rule{\linewidth}{0.6pt}
}
\newcommand{\examfooter}{%
  \vfill
  \rule{\linewidth}{0.4pt}\\[1pt]
  {\scriptsize\textit{Transcripci\'on digital --- MathUNAL}\hfill\textcolor{unalblue}{\textbf{\thepage}}}
}
\pagestyle{empty}

\begin{document}

\examheader

\textbf{Instrucciones:} La duraci\'on del examen es de 1 hora y 50 minutos. El examen consta de nueve preguntas en dos hojas impresas por ambos lados. Antes de empezar verifique que su examen est\'e completo y que sus datos sean correctos. No desengrape su examen ni a\~nada otras grapas. Utilice los espacios asignados para resolver cada pregunta. No est\'a permitido: hacer preguntas, usar calculadoras, sacar hojas en blanco, ni tomar apuntes en hojas diferentes a las del examen. No se permite que ning\'un celular se encuentre a la vista.

\bigskip

\textit{En las preguntas 1 a 5 indique si la afirmaci\'on dada es verdadera $\boxed{V}$ o falsa $\boxed{F}$, escoja la opci\'on correcta, o complete la respuesta en el espacio en blanco, seg\'un sea el caso. En estas preguntas la calificaci\'on tendr\'a en cuenta s\'olo la respuesta.}

\bigskip
\noindent\textbf{1.} (5\,\%) Si $D_{\vec u}f(x,y)$ existe, entonces
$$D_{-\vec u}f(x,y) = -D_{\vec u}f(x,y).$$
\hspace{2cm}$\boxed{\phantom{V}}$ V \qquad $\boxed{\phantom{V}}$ F

\bigskip
\noindent\textbf{2.} (5\,\%) Si
$$\lim_{x\to 0} f(x,0) = 0,$$
entonces
$$\lim_{(x,y)\to(0,0)} f(x,y) = 0.$$
\hspace{2cm}$\boxed{\phantom{V}}$ V \qquad $\boxed{\phantom{V}}$ F

\bigskip
\noindent\textbf{3.} (5\,\%) Sea
$$
f(x,y)=
\begin{cases}
x^2+xy^3, & y\geq 0,\\
0, & y<0.
\end{cases}
$$
Entonces la funci\'on $f$ es continua en el punto $(0,3)$ y es diferenciable en el punto $(3,5)$.

\hspace{2cm}$\boxed{\phantom{V}}$ V \qquad $\boxed{\phantom{V}}$ F

\bigskip
\noindent\textbf{4.} (5\,\%) Sea $f(x,y)=x^3-3x-y^3+3y$. Entonces

\begin{opciones}
  \item $f$ tiene un m\'inimo local (o relativo) en $(1,-1)$.
  \item $f$ tiene un m\'aximo local en $(1,-1)$.
  \item $f$ tiene un punto de silla en $(1,-1)$.
  \item No podemos afirmar que $f$ tiene un m\'aximo local o un m\'inimo local en $(1,-1)$.
\end{opciones}

\bigskip
\noindent\textbf{5.} (5\,\%) Sea $f=f(x,y)$ una funci\'on diferenciable con $f(1,2)=3$. Supongamos que la ecuaci\'on del plano tangente a la gr\'afica de $f$ en el punto $(1,2,3)$ es
$$z = 3+15(x-1)+5(y-2).$$
Si
$$\vec u = \frac{3}{5}\vec\imath - \frac{4}{5}\vec\jmath,$$
entonces
$$D_{\vec u}f(1,2) = \underline{\hspace{2.5cm}}.$$

\newpage
\examheader
\vspace{2pt}

\textbf{Preguntas con procedimiento}

\textit{En las preguntas 6 a 9 responda mostrando detalladamente el procedimiento utilizado.}

\bigskip
\noindent\textbf{6.} (20\,\%) Sea
$$
f(x,y)=
\begin{cases}
\dfrac{x^3-3y^3}{x^2+y^2}, & (x,y)\neq (0,0),\\[6pt]
0, & (x,y)=(0,0).
\end{cases}
$$
Encuentre $f_x(0,0)$ y $f_y(0,0)$ y determine si $f$ es diferenciable en $(0,0)$.

\vspace{4cm}

\noindent\textbf{7.} (20\,\%) Sea $f(x,y)=x^2-y^2$. Determine los vectores unitarios $\vec U$ para los cuales la derivada direccional de $f$ en el punto $\left(1,-\dfrac12\right)$ en direcci\'on de $\vec U$ tiene un valor de $1$.

\vfill
\examfooter
\newpage
\examheader
\vspace{2pt}

\noindent\textbf{8.} (20\,\%) Sea $f=f(x,y)$ una funci\'on diferenciable para la cual
$$\frac{\partial f}{\partial x} = \frac{y^2}{2} \qquad \text{y} \qquad \frac{\partial f}{\partial y} = xy.$$
Sea $w$ la funci\'on definida por
$$w(s,t) = f\left(st,\frac{s}{t}\right).$$
Encuentre el gradiente de $w$ en el punto $(2,1)$.

\bigskip
\bigskip

\noindent\textbf{9.} (15\,\%) Sea $f(x,y)=\sqrt{x^2+4y^2}$.
\begin{opciones}
  \item[(a)] Bosqueje las trazas de $f$ con los planos $xz$ y $yz$.
  \item[(b)] Determine los valores de $c$ para los cuales $f$ tiene una curva de nivel a altura $z=c$. Bosqueje las curvas de nivel de $f$.
  \item[(c)] Bosqueje la gr\'afica de $f$.
\end{opciones}

\vfill
\examfooter

\end{document}
